\section{Prime-carrying operators in the character sectors}
\label{sec:character-commutator}

The scalar character formula of MA-CH-001 has an actual operator
counterpart.  This section extends the one-sided construction
MA-MEY-001--003 to ideal Euler factors and checks the dual character
in the second diagonal.  The construction uses absolute Euler
convergence on its one-sided completions.  It does not use the
location of any completed zero.

\paragraph{MA-CC-001 (the character Euler commutator).}
Let $K$ be a number field and let $\eta$ be a primitive finite-order
Hecke character.  Choose a finite place set $S$ containing the
archimedean places, all ramification of $\eta$, and the finitely
many places needed for Meyer's structural local description and
the fixed additive character.  In the explicit imaginary quadratic
example of MA-CH-003, one may use $S_0$ from MA-CH-004.
Write $\mathscr I^S$ for the integral ideals prime to $S$,
$c_{\mathfrak a}=\eta(\mathfrak a)$, and $\mu_K$ for the ideal
M\"obius function.  For $b>1/2$, put
\[
\begin{aligned}
 \mathcal H_b&=L^2(\mathbb R,e^{2bx}\,dx),
       &T_\ell u(x)&=u(x+\ell),\\
 Z_{\eta,S}&=\sum_{\mathfrak a\in\mathscr I^S}
       c_{\mathfrak a}(N\mathfrak a)^{-1/2}
                         T_{\log N\mathfrak a},\\
 J_{\eta,S}&=\sum_{\mathfrak a\in\mathscr I^S}
       \mu_K(\mathfrak a)c_{\mathfrak a}
       (N\mathfrak a)^{-1/2}T_{\log N\mathfrak a}.
\end{aligned}
\tag{MA-CC.1}
\]
These are mutually inverse bounded operators, with absolutely
convergent operator series.  If $g\in C_c^\infty(\mathbb R)$ and
$\rho$ is smooth with compactly supported derivative and limits
$\rho_\pm$ at the two ends, the actual operator
\[
 K_{\eta,S,\rho}(g)
   =C_g Z_{\eta,S}[J_{\eta,S},M_\rho]
\tag{MA-CC.2}
\]
is trace class, and
\[
 \operatorname{Tr}K_{\eta,S,\rho}(g)
  =(\rho_- -\rho_+)
       \sum_{\substack{\mathfrak p\notin S\\k\ge1}}
       (\log N\mathfrak p)c_{\mathfrak p}^{\,k}
       (N\mathfrak p)^{-k/2}g(k\log N\mathfrak p).
\tag{MA-CC.3}
\]
The sum is finite.  Trace-norm estimates are uniform for a bounded
test family with fixed compact support.

Fix $\phi$ with $\phi_-=1$, $\phi_+=0$ and
$\phi(x)+\phi(-x)=1$.  In the actual compact-character sector
of Meyer's common space, the centered diagonal entries of
$P_+^S-P_-$ are
\[
 D_{\eta,S}=Z_{\eta,S}[J_{\eta,S},M_\phi],
 \qquad R D_{\bar\eta,S}R,
 \qquad Ru(x)=u(-x).
\tag{MA-CC.4}
\]
In particular the second term involves the dual Euler coefficients.
With $\beta(p_v)^{-1}=c_{\mathfrak p_v}$ as in MT-M22 and
probability Haar on $C_K^1$, set
\[
\begin{aligned}
 f_g(z)&=\beta(z)^{-1}|z|^{-1/2}g(\log|z|),\\
 B_{\eta,\overline S}(g)&=
 \sum_{\substack{\mathfrak p\notin S\\k\ge1}}
 \frac{\log N\mathfrak p}{(N\mathfrak p)^{k/2}}
 \left(c_{\mathfrak p}^{\,k}g(k\log N\mathfrak p)
 +\overline{c_{\mathfrak p}}^{\,k}g(-k\log N\mathfrak p)\right).
\end{aligned}
\tag{MA-CC.5}
\]
Then the common-space defect, computed on its actual nuclear
extension, satisfies
\[
 \mathcal D_S(f_g)
 =\operatorname{Tr}\lambda(f_g)(P_- -P_+^S)
 =-B_{\eta,\overline S}(g).
\tag{MA-CC.6}
\]
Consequently the operator identity in MA-HM-002 becomes
\[
 \operatorname{Tr}_{Z}C_g^Z
 =\chi_{\pi_+}(f_g)-\mathcal A_S(f_g)
                           -B_{\eta,\overline S}(g)
\tag{MA-CC.7}
\]
whenever its exact uniform-interior hypothesis holds.  The smooth
Meyer quotient identity before transport to $Z$ is unconditional.

\begin{proof}
Unique factorization of nonzero ideals replaces integer
factorization in MA-MEY-001.  Since $|c_{\mathfrak a}|=1$ outside
$S$ and $\|T_{\log N\mathfrak a}\|_{\mathcal H_b}
=(N\mathfrak a)^{-b}$, both norm series are dominated by
$\zeta_K(b+1/2)$.  Their product is absolutely convergent, and
the coefficient of the shift indexed by $\mathfrak c$ is
\[
 c_{\mathfrak c}(N\mathfrak c)^{-1/2}
                \sum_{\mathfrak d\mid\mathfrak c}\mu_K(\mathfrak d).
\]
This is zero unless $\mathfrak c=\mathcal O_K$, where it is one.
Distinct ideals may have the same norm; absolute convergence
allows this grouping by ideals before combining equal shifts.
Thus both operator products are the identity.

The compact-kernel estimate (MA-MEY.5), whose full proof is already
given, applies to each real shift.  It bounds the trace norm of
the term indexed by $(\mathfrak a,\mathfrak d)$ in (MA-CC.2) by
a constant times
\[
 (N\mathfrak a)^{-b-1/2}(N\mathfrak d)^{-b-1/2}
                                      (1+\log N\mathfrak d).
\]
The sum is finite.  Indeed for $\sigma=b+1/2>1$, choose
$0<\varepsilon<\sigma-1$ and use
$1+\log t\ll_\varepsilon t^\varepsilon$ together with the
absolute convergence of $\zeta_K(\sigma-\varepsilon)$.
The dependence on $g$ is through finitely many seminorms on its
fixed support, exactly as in (MA-MEY.5)--(MA-MEY.7).

Put $v=\log N\mathfrak a$ and $w=\log N\mathfrak d$.
The kernel of $C_gT_v[T_w,M_\rho]$ is
$g(x+v+w-z)(\rho(z)-\rho(z-w))$.  The weight conjugation is one
on the diagonal, so its trace is
$g(v+w)w(\rho_+-\rho_-)$ by (MA-MEY.9).
The absolute trace-norm convergence just proved permits summation
and regrouping.  The resulting coefficient at $\mathfrak c$ is
\[
 (\rho_+-\rho_-)
 c_{\mathfrak c}(N\mathfrak c)^{-1/2}g(\log N\mathfrak c)
 \sum_{\mathfrak d\mid\mathfrak c}
                          \mu_K(\mathfrak d)\log N\mathfrak d.
\]
The last sum is $-\Lambda_K(\mathfrak c)$, where
$\Lambda_K(\mathfrak p^k)=\log N\mathfrak p$ and it is zero
otherwise.  The proof is the same finite divisor expansion as
for integers: the coefficient of any one $\log N\mathfrak p$
is canceled by $\prod_{\mathfrak q\mid\mathfrak c,
\mathfrak q\ne\mathfrak p}(1-1)$ unless there is only one prime
divisor.  This proves (MA-CC.3).  There are only finitely many
integral ideals of bounded norm, proving finiteness of the trace sum.

It remains to check that these operators are the diagonal entries
of the actual adelic construction, including the dual orientation.
Write a raw vector in the $\beta$ sector as
\[
 v(ca(x))=\beta(c)^{-1}e^{-x/2}u(x),
 \qquad c\in C_K^1,\quad |a(x)|=e^x.
\tag{MA-CC.8}
\]
A finite-order character is trivial on this connected norm
splitting.  The raw translation $\lambda_{p_v}^{-1}$ becomes
$c_{\mathfrak p_v}(N\mathfrak p_v)^{-1/2}
T_{\log N\mathfrak p_v}$.  This follows either directly from
(MA-CC.8) or from (M22.4) with
$\beta(p_v)^{-1}=c_{\mathfrak p_v}$.
Thus the restriction of Meyer's $\mathcal L_S$ is exactly
$Z_{\eta,S}$, and its inverse is $J_{\eta,S}$.

The raw inversion $J_0v(z)=|z|^{-1}v(z^{-1})$ sends
(MA-CC.8) to
\[
 (J_0v)(ca(x))=\beta(c)e^{-x/2}u(-x).
\tag{MA-CC.9}
\]
It therefore exchanges the $\beta$ and $\beta^{-1}$ sectors and
becomes the linear reflection $R$ from $\mathcal H_{-b}$ to
$\mathcal H_b$ (or in the reverse direction).  It is not complex
conjugation of the vector.  The middle Euler multiplier, now in
the inverse character sector, has coefficients
$\overline{c_{\mathfrak p}}$ by unitarity.  Substitution in the
two explicit projection matrices (MA-MEY.13)--(MA-MEY.14)
proves (MA-CC.4), using $RM_{1-\phi}R=M_\phi$.

Meyer's Lemma~5.6 gives the nuclear extension of the complete
smoothed defect, including its off-diagonal entries, from the
paired weighted Hilbert spaces into the common smooth space.
Restriction to the fixed compact character and composition with
the continuous Hilbert inclusion give a trace-class operator.
The same nuclear rank-one representation computes its trace in
both spaces: each pairing is unchanged by that composition and
the scalar series is absolutely convergent.  Hence the trace is
the sum of the two Hilbert diagonal traces.  Since
$RC_gR=C_{g^\vee}$, (MA-CC.3) yields
\[
 \operatorname{Tr}\lambda(f_g)(P_+^S-P_-)
 =\operatorname{Tr}_{\mathcal H_b}C_gD_{\eta,S}
  +\operatorname{Tr}_{\mathcal H_b}C_{g^\vee}D_{\bar\eta,S}
 =B_{\eta,\overline S}(g).
\]
This proves (MA-CC.6).  Meyer's actual trace additivity
(MA-HM.23), and then the actual nuclear comparison (MA-HM.16),
give (MA-CC.7).  The first step uses no zero location; the
second uses precisely (MA-HM.2).
\end{proof}

\paragraph{MA-CC-002 (the non-Dirichlet cubic gap is an operator trace).}
For $K=\mathbb Q(\zeta_3)$, the character $\psi$ of
$K(269^{1/3})/K$ and the admissible set $S_0$ of MA-CH-003--004,
let $h$ be the test constructed in MA-CH-004 and put
$g(x)=h(e^x)$.  Then the actual one-sided prime-carrying operator
satisfies
\[
 \operatorname{Tr}K_{\psi,S_0,\phi}(g)=\log2,
 \qquad
 \operatorname{Tr}K_{\bar\psi,S_0,\phi}(g^\vee)=0,
 \qquad \mathcal D_{S_0}(f_g)=-\log2.
\tag{MA-CC.10}
\]
The genuine smooth quotient trace is consequently $-\log2$.
Under the same Hilbert interior-strip premise, the actual
orthogonal geometric cutoff limit is $+\log2$.

\begin{proof}
The previously checked construction has support strictly above
one, $h(4)=1$, and no other prime-ideal-power norm in its support.
The only contributing ideal is the unique prime above $2$,
of norm $4$, with $c_{\mathfrak p}=1$.
Its term in (MA-CC.3) is
$(\log4)/\sqrt4=\log2$; the reciprocal test has no positive
norm contribution.  This proves the three operator equalities.
MA-CH-004 also gives $A_{\mathbb C}(h)=0$ and $B_{\psi,S_0}(h)=0$.
The local Fourier-cutoff trace $\mathcal A_{S_0}(f_g)$ is
therefore zero by the local calculation, and the polar quotient
vanishes in this nontrivial character sector.  Trace additivity
gives the smooth quotient trace without comparing zero sums.
MA-HM-002 and MA-LIFT-003 give the last conclusion, with their
stated premise retained.  On wide strips it is unconditional;
on every positive narrow strip, RH for this particular character
is sufficient and is implied by $M$, not by the Dirichlet scope
of $A$.
\end{proof}

For primitive Dirichlet characters the same formulas use integer
ideals and the conjugate values $\overline\chi(p)^k$ on the
reciprocal terms.  The local real factor retains the parity in
(MA-CH.2).  The changes for an imaginary quadratic field are the
ideal norm, the complex local factor and conductor, and the dual
finite character in the reflected Euler multiplier.  The
operator proof itself needs no positivity and no simplicity of
zeros.  A compact cutoff has equal endpoint limits and hence
zero commutator trace in (MA-CC.3), so MA-MEY-002's obstruction
to replacing a one-sided cutoff by compact ones persists in
any character sector with a nonzero prime sum.

In the language of the requested geometric programme, the finite
character is a rank-one finite-monodromy local system on the
appropriate arithmetic open scheme, or its Artin motive over
$K$; duality replaces its Frobenius eigenvalues by their inverses.
These are the SGA~1, Expos\'e~V, \S7 and Deligne
\emph{Valeurs de fonctions $L$}, Definition~6.1 and
\S\S6.2--6.3 placements already used in MA-CH-003 and MA-S5.
The analytic commutator is an additional construction on the
idele-class representation, not an asserted proper algebraic
pushforward.  Its trace/extension mechanism has the inherited
cote~1, scans~079--080 and 111--112 comparison through
BR-TR-01/02 and F-SM-11/12.  None of these algebraic descriptions
removes the Hilbert strip-containment premise.

\paragraph{Exact source scope.}
The source maps are Meyer (2005), \path{arXiv:math/0311468v3},
Lemmas~5.3--5.6 and Theorem~5.8, author pp.~39--44; the
character and dual conventions are GRH\_NIGHT MT-M22,
(M22.1), (M22.4)--(M22.6).  The convergence and kernel estimates
used from this run are MA-MEY-001--003.  All ideal regroupings,
the reflected character computation and the operator specialization
are proved above.  The underlying algebraic character and its
local factors are the already proved MA-CH-003--004.
